\section{Introduction}

\subsection{Real-world Context}
Urban Traffic Congestion is a pervasive issues across a range of different demographics, whether it is in a developing or developed country. 
Traffic Congestion has always been a significant concern in traffic control mechanisms caused by an unbalanced demand-supply situation that directly affects the problem of large fuel consumption and greenhouse emission.
Congestion leads can lead to a range of issues such as traffic jams, unhealthy air quality, vibrations, noises and accidents which are significant problems for the environment and human health \cite{jilaniSystematicReviewUrban2025}. 

In the UK, this issue is particularly prevalent. 
In the 2019, UK drivers lost an average of 115 hours each to road traffic congestion \cite{pishue2020INRIXGlobal}.

To reduce congestion, we should look at optimising traffic flow, this can be done through the use of traffic lights, by employing various strategies we can maximise the flow of traffic and reduce congestion. Since traffic lights are already widely used, it is important to understand how they can be optimised to improve traffic flow and reduce congestion allowing for minimal costs.

\subsection{Why discrete cellular automata}
We can leverage the use of cellular automata to model traffic flow, this is because they are a powerful tool for simulating complex systems and can capture the dynamics of traffic flow in a simple and efficient way. Cellular Automata also allow us to assess the impact of different traffic light strategies on traffic flow, and can help us to identify the most effective strategies for reducing congestion.

Traffic cellular automata models have been widely used in traffic flow research, and have been shown to be effective in capturing the dynamics of traffic flow and congestion. Examples of cellular automata models include the Nagel-Schreckenberg (NaSch) model and the Biham-Middleton-Levine (BML) model \cite{nagelCellularAutomatonModel1992,bihamSelforganizationDynamicalTransition1992}.

For this study, we will use the NaSch model due to its simplicity and more clearly explaining the effects on traffic lights in the system. With this model we can easily consider the effects of multiple traffic lights on the system as well as the effects of different traffic light strategies on traffic flow and congestion. Additionally we can also consider how a second lane can affect the traffic flow and congestion in the system with different vehicle types taken into consideration.


\subsection{Aims and Research Questions}

In this report, we will explore three different contexts of traffic flow, exploring how different traffic light layouts effect the system. 
The first context will be a single lane system, where we vary traffic light number and traffic light strategies \cite{nagelCellularAutomatonModel1992}.
The second one will be a two lane system, where will explore the effect of varying traffic light number and traffic light strategies with the introduction of different vehicle types on the system \cite{rickertTwoLaneTraffic1996,chowdhuryParticleHoppingModels1997}. 
The last one is a perpendicular intersection in which two single lane roads intersect with each other, we will explore the effect of varying traffic light strategies on the system \cite{foulaadvandVehicularTrafficFlow2011}.
The unifying question across these two is light configuration and strategy, we want to understand how different traffic light configurations and strategies can affect traffic flow and congestion in both single lane and two lane systems. The three geometries are shown in Figure~\ref{fig:scenarios_schematic}.

We will address four research questions:
\begin{enumerate}
    \item How does single-lane through-flow scale with the number of traffic lights, the cycle length, the green fraction and the inter-light phase offset?
    \item When does green-wave coordinate outperform synchronised lights and why?
    \item Does symmetric lane changing complement the green-wave or fight it?
    \item How do the above results change when we consider a perpendicular intersection?
\end{enumerate}




